\subsection{根}

G 相对于 T 的根是 G 的伴随表示的非零权。G 相对于 T 的根所成的集记作 R(G, T)，在不致混淆时简记为 R。由命题 6，映射

\[
\delta : \mathrm{X(T)} \to \mathfrak {t} _ {\mathbf {C}} ^ {*}
\]

\((t_{\mathbf{C}}^{*} \text{ 表示复向量空间的对偶 } t_{\mathbf{C}}) \text{ 将 } R(G, T) \text{ 双射地映到集合 } R(\mathfrak{g}_{\mathbf{C}}, t_{\mathbf{C}}) \text{ 分裂约化代数的根集 } (\mathfrak{g}_{\mathbf{C}}, t_{\mathbf{C}}) \text{ (Chap. VIII, §2, no. 2, Remark 4). 若置}\)

\[
\mathfrak {g} ^ {\alpha} = (\mathfrak {g} _ {\mathbf {C}}) _ {\alpha} (\mathrm{T}) = (\mathfrak {g} _ {\mathbf {C}}) _ {\delta (\alpha)} (\mathfrak {t} _ {\mathbf {C}}),\tag{11}
\]

则对一切 \(\alpha\in R\) ，每个 \(\mathfrak{g}^{\alpha}\) 在 C 上维数为 1（同前，定理 1），且

\[
\mathfrak {g} _ {\mathbf {C}} = t _ {\mathbf {C}} \oplus \bigoplus_ {\alpha \in R} \mathfrak {g} ^ {\alpha}.\tag{12}
\]

对每个 \(\alpha \in \mathbb{R}\) ，用 \(V(\alpha)\) 表示 \(\mathfrak{g}\) 的 2 维子空间，满足 \(V(\alpha)_{(\mathbf{C})} = \mathfrak{g}^{\alpha} + \mathfrak{g}^{-\alpha}\) ；\(\mathfrak{g}\) 对 T 的伴随表示的非零同构型分支是 \(\mathfrak{t}\) 与诸 \(V(\alpha)\) 。进而，设 K 是与 \(\mathfrak{g}\) 的 Killing 型相伴的二次型；它是负的（§1 第 3 节，命题 1），而它的限制 \(\mathrm{K}(\alpha)\) （到 \(\mathrm{V}(\alpha)\) 上的）是负的且分离的。对 T 的每个元素 t，Ad t 保持 \(\mathrm{K}(\alpha)\) 不变，从而给出李群态射

\[
\iota_ {\alpha}: \mathrm{T} \to \mathbf {S O} (\mathrm{K} (\alpha)).
\]

存在唯一的同构 \(\rho_{\alpha}:\mathbf{U}\to \mathbf{SO}(\mathrm{K}(\alpha))\) 使得 \(\iota_{\alpha} = \rho_{\alpha}\circ \alpha\) 。事实上，设 \(X\) 是 \(\mathfrak{g}^{\alpha}\) 的非零元素，\(Y\) 是 \(X\) 在 \(\mathfrak{g}_{\mathbf{C}}\) 相对于 \(\mathfrak{g}\) 的共轭下的像；则 \(Y\in \mathfrak{g}^{-\alpha}\) ，且得到基 \((U,V)\) （\(\mathrm{V}(\alpha)\) 的），其取法为 \(U = X + Y,V = i(X - Y)\) ；\(\mathrm{V}(\alpha)\) 的自同态（由 \(\operatorname {Ad}t,t\in \mathrm{T}\) 诱导的）关于基 \((U,V)\) 的矩阵是

\[
\left( \begin{array}{c c} \mathcal {R} (t ^ {\alpha}) & - \mathcal {I} (t ^ {\alpha}) \\ \mathcal {I} (t ^ {\alpha}) & \mathcal {R} (t ^ {\alpha}) \end{array} \right),
\]

断言得证。

命题 8. 设 \(\mathrm{Q}(\mathrm{R})\) 是 \(\mathrm{X}(\mathrm{T})\) 的、由 \(\mathrm{G}\) 的诸根生成的子群。

\begin{enumerate}
\def\labelenumi{\alph{enumi})}
\item
  G 的中心 \(\mathrm{C}(\mathrm{G})\) 是 T 的闭子群，等于诸根之核的交。典则映射 \(\mathrm{X}(\mathrm{T}/\mathrm{C}(\mathrm{G})) \to \mathrm{X}(\mathrm{T})\) 单，且其像为 \(\mathrm{Q}(\mathrm{R})\) 。
\item
  紧群 C(G) 同构于离散群 X(T)/Q(R) 的对偶（《谱理论》，第 II 章，§1，第 1 节，定义 2）。
\item
  C(G) 约化为单位元当且仅当 Q(R) 等于 X(T)。
\end{enumerate}

由 §2 第 2 节定理 2 的推论 2，C(G) 含于 T。由于它是伴随表示的核，它是诸根之核的交，换言之是 X(T) 的子群 Q(R) 的正交补。于是，命题由《谱理论》第 II 章 §1 第 7 节定理 4 与第 5 节定理 2 得出。

命题 9. 在 T 上诱导恒等映射的 G 的每个李群自同构都形如 Int t，其中 \(t \in T\) 。

先假定 \(\mathrm{C}(\mathrm{G})\) 约化为单位元，换言之 \(\mathrm{X}(\mathrm{T}) = \mathrm{Q}(\mathrm{R})\) （命题 8）。设 f 是在 T 上诱导恒等映射的 G 的自同构，\(\varphi = \mathrm{L}(f)_{(\mathbf{C})}\) ；则 \(\varphi\) 是 \(g_{C}\) 的自同构且在 \(t_{C}\) 上诱导恒等映射。由第 VIII 章 §5 第 2 节命题 2，存在唯一的同态 \(\theta : \mathrm{Q}(\mathrm{R}) \to \mathbf{C}^{*}\) ，使得 \(\varphi\) 在每个 \(\mathfrak{g}^{\alpha}\) 上诱导比为 \(\theta(\alpha)\) 的位似。由于 \(\varphi\) 保持 \(g_{C}\) 的实形式 g 稳定，它与共轭 \(\sigma\) （\(g_{C}\) 相对于 g 的）交换；而 \(\sigma(\mathfrak{g}^{\alpha}) = \mathfrak{g}^{-\alpha}\) ，故 \(\theta(-\alpha) = \overline{\theta(\alpha)}\) 对一切 \(\alpha \in R\) 成立。这蕴含 \(\theta(\alpha)\overline{\theta(\alpha)} = \theta(\alpha)\theta(-\alpha) = 1\) 。由此得出 \(\theta\) 取值于 U，从而由对偶性对应 T 的元素 t，使得 \((\operatorname{Ad} t)_{(\mathbf{C})} = \varphi\) ，故 \(\operatorname{Int} t = f\) 。

在一般情形，把上述应用于其中心约化为单位元的群 G/C(G) 及其极大环面 T/C(G)。

由此得出，若 \(f\) 是 \(G\) 的、在 \(T\) 上诱导恒等映射的自同构，则存在元素 \(t\) （\(T\) 中的），使 \(f\) 与 \(\operatorname{Int} t\) 过渡到商后在 \(G / C(G)\) 上诱导相同的自同构。而由于典则态射 \(D(G) \to G / C(G)\) 是有限覆盖（§1 第 4 节，命题 4 的推论 1），\(f\) 与 \(\operatorname{Int} t\) 在 \(D(G)\) 上、从而在 \(D(G) \times C(G)\) 上、进而也在 \(G\) 上（同前）诱导相同的自同构。

推论. 设 u 是 G 的自同构，H 是 G 的由 u 的不动点组成的闭子群。则自同构 u 是内的当且仅当 \(H_{0}\) 极大秩。

若 u 等于 Int g（\(g \in G\) ），则子群 \(H_{0} = \mathrm{Z}(g)_{0}\) 极大秩（ \(\S2\) ，第 2 节，推论 3）。反之，若 H 包含一个极大环面 S，则自同构 u 形如 Int s，\(s \in S\) （命题 9）。
% semantic-fragments: legacy-input-seam
\relax{}
